\fancyhead[L]{\small MAGNETISCHE BEOBACHTUNGEN.}
\SourceLeaf{523}\pdfbookmark[0]{Nordlicht vom 7. Januar 1831}{nordlicht_1831_januar}
\begin{center}
\hrule\vspace{2mm}
Göttingische gelehrte Anzeigen. 1831. Februar 28.
\vspace{2mm}\hrule\vspace{5mm}
\end{center}
Professor \textsc{Gerling} in Marburg hat der Königlichen Sozietät eine Notiz über seine Beobachtung
\begin{center}\itshape des am 7. Januar dieses Jahres gesehenen Nordlichts\end{center}
\noindent vorgelegt, die im Allgemeinen mit den bereits von anderen Orten bekannt gewordenen Beobachtungen übereinstimmt. Sie berührt aber auch einen Umstand, der besondere Aufmerksamkeit verdient und bisher, wie es scheint, nicht genügend gewürdigt worden ist. Deshalb teilen wir hier einen Auszug daraus mit.

Das Phänomen wurde in Marburg schon ab 6 Uhr beobachtet. Herr \textsc{Gerling} erfuhr jedoch erst um 8 Uhr davon. Zu diesem Zeitpunkt war am gesamten nördlichen Himmel bis zum tiefsten Punkt, den er aus den Fenstern seiner Wohnung sehen konnte, nichts Ungewöhnliches zu erkennen. Gegen 9 Uhr zeigten sich aber wieder auffallende rote Streifen am nördlichen Himmel. Herr \textsc{Gerling} ging sofort auf den Schlossberg, von dem aus man freie Sicht hat, um noch so viel wie möglich von der Erscheinung wahrzunehmen.

Zunächst waren über eine Ausdehnung von etwa 50—60 Grad zwischen Nordosten und Nordwesten nur rote Streifen und Flecken am Himmel zu sehen. Sie erstreckten sich ohne durchgehenden Zusammenhang über den angegebenen Bogen im Azimut und erreichten im Mittel etwa 45 Grad Höhe. In der Mitte dieses Azimutbo\SourceLeaf{524}gens, um den Meridian herum, zeigten sich am sonst klaren Himmel über eine geschätzte Ausdehnung von etwa 30—40 Grad im Azimut schwarze Flecken. Ihrem Aussehen nach waren sie nur mit kleinen schwarzen Wolken zu vergleichen. Diese Flecken vermehrten sich allmählich und flossen schließlich zu dem dunklen Segment zusammen, das nach allen Beschreibungen für das Nordlicht charakteristisch zu sein scheint. Zugleich nahmen die zuerst erwähnten roten Flecken an Intensität zu und ordneten sich strahlenförmig auf das schwarze Segment hin an. Von diesem gingen dann zwischen den roten Strahlen auch weiße und gelbliche aus, die ohne auffallend plötzliches Emporschießen etwa 50 Grad Höhe erreicht haben mochten.

‘Bis hierher, fährt Herr \textsc{Gerling} fort, scheint diese Beobachtung mit dem, was andere Beobachter zur gleichen Zeit und bei früheren Nordlichtern gesehen haben, ganz übereinzustimmen. Sie wäre daher kaum erwähnenswert, wenn mir dabei kein Umstand aufgefallen wäre, der meines Wissens weder bei diesem Nordlicht noch, soweit ich feststellen konnte, sonst zur Sprache gekommen ist. Die Sterne des Schwans, über die die intensiven weißen und roten Strahlen hinweggingen, und der Stern $\alpha$ in der Leier, \textit{der tief im schwarzen Segment stand}, verloren augenscheinlich überhaupt nichts an Sichtbarkeit und scheinbarer Helligkeit. Diese Tatsache scheint für die rätselhafte Frage, was es mit dem dunklen Segment eigentlich auf sich hat, wenigstens das negative Ergebnis zu liefern, dass es \textit{keine gewöhnliche Wolke} ist. Eine solche könnte das Sternlicht nämlich nicht durchlassen.’

Schon beim Nordlicht vom 22. Oktober 1804 bemerkte \textsc{Wrede}, ohne allerdings diesen Grund anzuführen, dass man das dunkle Segment \textit{zu Unrecht} eine Wolke nenne. \textsc{Gilbert} verteidigt dagegen den Ausdruck und fügt hinzu, er habe im dunklen Segment nichts bemerkt, was ihn auf den Gedanken hätte bringen können, dort etwas anderes als eine dunkle Wolke zu sehen. Auch \textsc{Mayers} Ansicht im Handbuch der physischen Astronomie, dass die dichtere, mit Dünsten erfüllte Luft am Horizont ausreiche, das dunkle Segment zu erklären, scheint sich mit der von Herrn \textsc{Gerling} beobachteten Tatsache nicht vereinbaren zu lassen.

Herr \textsc{Gerling} fügt hinzu, dass in den früheren Stunden, als sich das in seiner Ausdehnung veränderliche Segment sehr hoch erstreckte, ein glaubwür\SourceLeaf{525}diger Zeuge den Stern $\alpha$ der Leier \textit{im Segment} so hell wie zu jeder anderen Zeit habe glänzen sehen. Ein anderer habe zu einer Zeit, als das dunkle Segment noch nicht bis zu diesem Stern reichte, \textit{andere Sterne im Segment} erblickt.

Herr \textsc{Gerling} hat außerdem einen Auszug aus seinem meteorologischen Journal vom 5.—9. Januar beigefügt. Er enthält jedoch außer einem Anstieg des Barometers um drei Viertel Zoll vom Nachmittag des 6. Januar bis zum Abend des 7. Januar nichts Auffallendes. Der Wind kam am 7. Januar aus Norden.

Die Beobachtungen, die Professor \textsc{Harding} hier in Göttingen an diesem Nordlicht gemacht hat, stimmen im Wesentlichen mit den von anderen Orten bekannt gewordenen überein. Erwähnenswert ist jedoch, dass sich die Magnetnadel während des Phänomens von ihrer gewöhnlichen Stellung um etwa drei Viertel Grad nach Norden bewegte und am nächsten Morgen wieder in diese Stellung zurückgekehrt war.

\bigskip
